\documentclass[11pt]{article}
\usepackage[T1]{fontenc}
\usepackage[utf8]{inputenc}
\usepackage{lmodern,amsmath,amssymb,amsthm}
\usepackage[margin=1in]{geometry}
\usepackage[hidelinks]{hyperref}
\newtheorem{theorem}{Theorem}
\newtheorem{lemma}[theorem]{Lemma}
\newcommand{\Z}{\mathbb Z}\newcommand{\N}{\mathbb N}
\newcommand{\C}{\mathbb C}\newcommand{\E}{\mathbb E}
\newcommand{\Symb}{\operatorname{Symb}}
\newcommand{\Nil}{\operatorname{Nil}}
\newcommand{\one}{\mathbf1}
\title{The Section 13 correlation and symmetry join}
\author{Author proposal: /root/astra\_v14\_baseline\_compiler}
\date{7 November 2026}
\begin{document}\maketitle
\noindent\textbf{Status.} This is a separate author proposal following
LP05-001. It does not alter that frozen proof or claim a complete inverse
theorem. Its project premises and unchanged published inputs are explicit
below. Mathematical acceptance, related-use completeness and incorporation
require independent decisions. The signed identities here retain the actual
unconjugated $\alpha$ in published (13.10).

\section{Target and exact inputs}
Fix a standard integer $s\geq3$ and an unbounded limit integer $N$.
Write $I_t=[-tN,tN]\cap{}^*\Z$, $J=[N]$, $r=s-1$ and $m=s-2$.
All geometric data and finite dimensions are standard. Vertical identities
hold for every limit point and every limit top-group element under the
natural ultrapower extension, as in the project tautification and LP13
units. Tensor products have their Hilbert tensor norms. Every nilcharacter
is unit-vector-valued; a scalar coordinate need not have modulus one.

\begin{theorem}[Theorem 7.4 on the consumed domain]\label{thm:target}
Let $f$ be bounded on $J$, extended by zero outside $J$, and let
$H\subseteq I_1$ be a dense limit subset. Suppose that $\chi(h,n)$ is
a globally represented multidegree-$(1,s-1)$ nilcharacter and
$\Delta_h f$ has degree-at-most-$(s-2)$ correlation with $\chi(h,\cdot)$
on $J$ for every $h\in H$, with GTZ's correlation convention.
Then there are a globally represented degree-at-most-$s$ unit nilcharacter
$\Theta$ and a bounded globally represented nilsequence $\Psi(h,n)$
supported on
\[
 \mathcal J=\{(i,j):i+j\leq s-1\}\cup\{(i,s-i):2\leq i\leq s\}
                                                        \tag{1}
\]
such that, on $I_1\times J$, each component of $\chi(h,n)$ is a bounded
finite linear combination of components of
\[
 \Theta(n+h)\otimes\overline{\Theta(n)}\otimes\Psi(h,n).  \tag{2}
\]
All families and bounds are standard in the limit sense, independently of
the interval length. No effective numerical constants are asserted.
\end{theorem}

We use the following exact premises.
The project LP05-001 supplies a unit nilcharacter $\widetilde\chi$ of
cube multidegree $(1,\ldots,1)$, literally symmetric in its last $r$
variables, with
\[
 [\chi(h,n)]_{\Symb_s({}^*\Z^2)}
       =s[\widetilde\chi(h,n,\ldots,n)]_{\Symb_s({}^*\Z^2)}. \tag{3}
\]
The project LP13 theorem supplies multilinearity of the symbol of
$\widetilde\chi$ under fixed integer linear substitutions, in each slot.
The project tautification-003 Theorem 3 supplies the uniform shift tensor
of a positive-degree nilcharacter, lowering degree by one in the ordinary
degree case. It applies to limit shifts and limit families.

The unchanged published inputs are the definitions and elementary symbol
arithmetic E.3, E.7 and E.8(i)--(iii),(vi); nilcharacter approximation and
selection E.5/E.6 and Appendix A; Proposition 8.4 with its plain-correlation
hypothesis; the converse Proposition 5.6; and E.11/E.13. Their exact source
applications are recorded separately. In particular Proposition 8.4 is a
distinct input. The weaker degree-$(s-2)$ biased-quadruple conclusion of
Proposition 7.3 alone would not retain the lower factors needed here.
No corrected assertion is supplied by citing the unpublished erratum.

For an ordered list of linear forms define
\[
 T(L_1;L_2,\ldots,L_s)=[\widetilde\chi(L_1,L_2,\ldots,L_s)]_{\Symb_s}.
\]
The semicolon separates the distinguished first slot. We use the notation
$n^j$ in a slot list for $j$ repeated entries, not a power of $n$.
All formulas between such symbols are on the full indicated limit lattice
until a box is explicitly named.

\section{Plain correlation and retention of the lower factors}
By E.5/E.6 in the ordinary one-dimensional degree ordering and Appendix A
selection, after restricting to a still dense
internal subset of $H$, choose degree-at-most-$(s-2)$ unit nilcharacters
$\varphi_h$ so that $\Delta_hf$ plainly correlates with
$\chi(h,\cdot)\otimes\varphi_h$. The family is internal, with a standard
finite geometry list, bounded dimensions and common Lipschitz bounds.
Here the selection is into the countable increasing complexity presentation
of the nilcharacter class. An internal image contained in that union is
contained in a standard stage by Appendix A compactness; otherwise nested
nonempty internal sets of parameters outside successive stages give an
image outside the union. Positive correlations can likewise be made
uniform by A.6. Pad targets by zero coordinates to one fixed dimension.
Outside the retained set use a fixed constant unit nilcharacter, and keep
all parameter indicators separately. This preserves the unit norm needed
for the identities below.
Conjugate the selected lower nilcharacter if needed to match the published
correlation convention. A finite pigeonhole also fixes a scalar coordinate
of $\chi(h,\cdot)\otimes\varphi_h$ retaining plain correlation on a dense
set. Thus the scalar version of Proposition 8.4 suffices; the full unit
vectors remain available behind these coordinates.

Apply published Proposition 8.4 to this actual family. Its hypotheses are
plain correlation on $J$, internal dependence, bounded functions and a
dense subset of $I_1$; all hold. Its conclusion gives bias for many additive
quadruples $(h_1,h_2,h_3,h_4)$ in that subset, with tensor signs $+,+,-,-$
and the shift $n+h_1-h_4$ on the second and fourth factors. Parameterise
\[
 (h_1,h_2,h_3,h_4)=(h_0+a,h_0+b,h_0+a+b,h_0).
\]
This is an injective correspondence with the retained quadruples.
$h_0\in I_1$ and $a,b\in I_2$. Pigeonhole in $h_0$ to obtain one fixed
$h_0$ and a dense set of pairs $(a,b)\in I_2^2$.
Define the global unit nilcharacter
\[
 \tau(a,b,n)=\chi(h_0+a,n)\otimes\chi(h_0+b,n+a)
 \otimes\overline{\chi(h_0+a+b,n)}
 \otimes\overline{\chi(h_0,n+a)}.                         \tag{4}
\]
It has total degree at most $s$, not degree three for general $s$.

For any coordinate of a unit nilcharacter $v$, the exact identity
\[
 v_j(n+a)=\sum_k\bigl(v_j(n+a)\overline{v_k(n)}\bigr)v_k(n) \tag{5}
\]
follows from $\sum_k|v_k(n)|^2=1$. The bracket is a coordinate of the
shift tensor. Apply the project shift theorem to $\varphi_{h_0+b}$ and
$\varphi_{h_0}$, with shift $a$. These brackets have degree at most
$s-3$, uniformly in all parameters. Conjugate (5) when the shifted factor
is conjugated. Thus the original bias implies degree-at-most-$(s-3)$ bias
of
\[
 \tau(a,b,n)\otimes\varphi_{h_0+a}(n)\otimes\varphi_{h_0+b}(n)
 \otimes\overline{\varphi_{h_0+a+b}(n)}
 \otimes\overline{\varphi_{h_0}(n)}.                      \tag{6}
\]
Selection gives internal scalar degree-at-most-$(s-3)$ tests with common
bounds and finite geometries. A standard finite choice of coordinates
retains a positive density of the pairs. Phase-align their scalar averages
by a bounded parameter scalar, and set the test equal to zero on other
pairs. These factors are constant in $n$, so they have degree zero.
We obtain fixed scalar coordinates $t$ of $\tau$ and $u,v,w,z$ of the
displayed $\varphi$ factors with
\[
 \left|\E_{a,b\in I_2,n\in J}
 t(a,b,n)u_a(n)v_b(n)\overline{w_{a+b}(n)}\overline{z(n)}
 p_{a,b}(n)\right|\gg1.                                 \tag{7}
\]
The scalar $p_{a,b}$ is a uniformly bounded internal lower-degree family.
Finite coordinate selection, rather than division by a coordinate, justifies
this passage even if individual coordinates vanish.

\section{The three Cauchy--Schwarz eliminations}
Every average in this section has $n\in J$. All bounds are standard.
Applying Cauchy--Schwarz in $(a,n)$ to remove the bounded factor
$u_a\overline z$ in (7), and expanding the squared inner average, gives
\[
 \left|\E_{a,b,b'\in I_2,n\in J}
 t(a,b,n)\overline{t(a,b',n)}v_b(n)\overline{v_{b'}(n)}
 \overline{w_{a+b}(n)}w_{a+b'}(n)
 p_{a,b}(n)\overline{p_{a,b'}(n)}\right|\gg1.             \tag{8}
\]
Put $c=a+b+b'$. Keep $b,b'\in I_2$, take $c\in I_6$, and retain
$\one_{c-b-b'\in I_2}$. This exact indicator is independent of $n$.
The normalized average changes by the factor $|I_2|/|I_6|$, comparable
to one. No three-term sum of elements of $I_2$ is asserted to lie in $I_3$.
Pigeonhole in $c$ to fix $c_0\in I_6$. Set
\[
 \alpha(b,b',n)=\tau(c_0-b-b',b,n)
                    \otimes\overline{\tau(c_0-b-b',b',n)},\qquad
 V_b(n)=v_b(n)w_{c_0-b}(n).                              \tag{9}
\]
In particular the second factor of $V_b$ is not conjugated: its sign
follows from the last two $w$ factors in (8). Absorb the parameter cutoff
into the scalar lower test. Equation (8) becomes, for a fixed coordinate
$a_0$ of the actual $\alpha$,
\[
 \left|\E_{b,b'\in I_2,n\in J}
 a_0(b,b',n)V_b(n)\overline{V_{b'}(n)}q_{b,b'}(n)\right|\gg1. \tag{10}
\]
The functions $V_b$ need only be bounded; no inverse or unit norm for them
is used. Cauchy--Schwarz in $(b',n)$ removes $\overline{V_{b'}}$ and
produces the product
$a_0(b_1,b',n)\overline{a_0(b_2,b',n)}V_{b_1}\overline{V_{b_2}}$
with the product of the corresponding $q$ tests. Cauchy--Schwarz in
$(b_1,b_2,n)$ removes the remaining $V$ factors. Its expanded four-term
product is
\[
 a_0(b_1,b'_1,n)\overline{a_0(b_2,b'_1,n)}
 \overline{a_0(b_1,b'_2,n)}a_0(b_2,b'_2,n).               \tag{11}
\]
Its scalar test is the product of four $q$ tests with the same signs,
still degree at most $s-3$, internal and uniformly bounded.
Pigeonhole in $(b_2,b'_2)$, calling the fixed values $(b_0,b'_0)\in I_2^2$.
The resulting expression is a coordinate of the global unit nilcharacter
\[
 \alpha'(b,b',n)=\alpha(b,b',n)\otimes\overline{\alpha(b_0,b',n)}
 \otimes\overline{\alpha(b,b'_0,n)}\otimes\alpha(b_0,b'_0,n). \tag{12}
\]
Hence $\alpha'(b,b',\cdot)$ is degree-at-most-$(s-3)$ biased on $J$
for many $(b,b')\in I_2^2$. To see the last assertion quantitatively,
if the final average has modulus at least $\delta>0$ and each inner
average is at most $B$, a proportion at least $\delta/(2B)$ has modulus
at least $\delta/2$. All selections and parameter cutoffs above depend
internally on the parameters. Products and conjugates do not increase the
ordinary nilsequence degree or destroy the fixed finite complexity bounds.

\section{Exact symbols and the lower-slot remainder}\label{sec:algebra}
Constants such as $h_0,c_0,b_0,b'_0$ are limit integers. Give the original
multidegree group its total-degree filtration generated by all components
whose total weight is \emph{at least} the index, as verified in LP05.
Apply the uniform project ordinary-degree shift theorem on the full slot
domain to a constant translation in the relevant coordinate. Its difference
has degree at most $s-1$, and stays lower-degree under any fixed integer
linear pullback. Thus a constant added to a slot can be removed in the
degree-$s$ symbol. A slot consisting only of a constant becomes zero.
LP13 additivity gives $T(\ldots,0,\ldots)=2T(\ldots,0,\ldots)$, so that
symbol is zero. This argument does not choose an implicit subgroup
basepoint for a fixed-slot orbit. The coarsening uses $|a|\geq j$;
the opposite inequality printed in Example 6.11 is not used.

From (3),(4), linearity in the first slot, and last-block symmetry,
\begin{align}
 [\tau(a,b,n)]
 &=s\bigl(T(b;(n+a)^r)-T(b;n^r)\bigr)\notag\\
 &=s\sum_{j=1}^r {r\choose j}T(b;a^j,n^{r-j}).           \tag{13}
\end{align}
In the first line $(n+a)^r$ denotes repeated slots.
The leading term is $srT(b;a,n^m)$. Let
\[
 R_\tau(a,b,n)=s\sum_{j=2}^r {r\choose j}T(b;a^j,n^{r-j}). \tag{14}
\]
Every summand has at most $s-3$ slots equal to $n$.
Substitute $a=c_0-b-b'$ in (9), remove $c_0$ as just justified, and
expand the $j$ copies of $-b-b'$ by LP13. The leading symbol of $\alpha$ is
\[
 -sr\,T(b-b';b+b',n^m).                                \tag{15}
\]
This has the order of the two distinguished entries dictated by (13).
Interchanging them would require an unproved first-slot symmetry and is
not permitted. The remainder is the difference of the two substitutions
in (14), each fully expanded into entries in $\{b,b',n\}$ and still with
at most $s-3$ copies of $n$.

Apply the mixed difference (12). At the level of degree-$s$ symbols,
constant slots vanish; a monomial with no $b$ or no $b'$ cancels.
In (15) the pure terms $T(b;b,n^m)$ and $T(b';b',n^m)$ therefore disappear.
There is a globally represented unit nilcharacter $\beta$ such that
\[
 [\alpha'(b,b',n)]=-sr\bigl(T(b;b',n^m)-T(b';b,n^m)\bigr)+[\beta(b,b',n)].
                                                               \tag{16}
\]
Here $[\beta]$ is an explicit finite integer linear combination of symbols
$T(v_1;\ldots,v_s)$, with every $v_i\in\{b,b',n\}$ and at most $s-3$
slots equal to $n$. More precisely, expand (14) after the two substitutions
$(a,b)=(-b-b',b)$ and $(-b-b',b')$, subtract the second, then delete the
monomials missing $b$ or missing $b'$. Tensor the corresponding actual
pullbacks of $\widetilde\chi$ with their signed integer multiplicities.
This constructs $\beta$ and proves (16). Thus the needed remainder
assertion is stronger than merely a separate bound on its degree in $n$.

Define the actual unit nilcharacter
\[
 \sigma(b,b',h_1,\ldots,h_m,n)=
 \bigotimes_{\omega\in\{0,1\}^m}\mathcal C^{|\omega|}
   \alpha'\left(b,b',n+\sum_{j=1}^m\omega_jh_j\right),    \tag{17}
\]
where $\mathcal C$ is conjugation. Every shift has its own $h_j$.
It has total degree at most $s$ by linear pullback and tensor arithmetic.
Apply the alternating sum to (16). In a term with fewer than $m$ slots
equal to $n$, expand each such slot into $n$ and the available $h_j$.
If a subset $S$ of the shift labels is used, its coefficient contains
\[
 \sum_{\omega:\,\omega_j=1\ (j\in S)}(-1)^{|\omega|}=0
 \quad\hbox{unless }S=[m].
\]
There are too few slots to use all $m$ labels, so every $\beta$ term
vanishes. For the leading term there are exactly $m!$ surviving bijections
of labels to slots, each with coefficient $(-1)^m$. Last-block symmetry
identifies their values. With the minus sign in (16), the exact result is
\[
 [\sigma]=(-1)^{s-1}s!
 \bigl(T(b;b',h_1,\ldots,h_m)-T(b';b,h_1,\ldots,h_m)\bigr). \tag{18}
\]
The unused coordinate $n$ on the right is still a coordinate of the same
domain. Formula (18) keeps the actual definitions (9),(12),(17). If one
instead conjugates $\alpha$ and $\alpha'$ explicitly, its coefficient
changes sign; silently using the other object is not an identification.

\section{From correlation to a zero symbol on the correct box}
By Proposition 5.6, the positive degree-at-most-$(s-3)$ bias obtained
after (12) implies a positive $U^m[J]$ norm for a scalar coordinate of
$\alpha'$, for many $(b,b')$. Appendix A uniformity gives a standard
positive lower bound on a dense internal set. Equivalently the average
of the vector $U^m[J]^{2^m}$ quantities is $\gg1$. Since the dimension is
standard finite, at least one fixed scalar coordinate contributes $\gg1$.
For $s=3$, this is the converse at degree zero and $U^1$, so no extra
one-step test has been inserted.

Expand the interval Gowers norm using its zero-extension definition and
the same normalization as GTZ. One can choose the auxiliary cyclic group
of order a fixed multiple of $N$ greater than $2(m+1)N$. For a nonzero
cube all vertices lie in $J$, so its differences have unique lifts
$h_j\in I_1$, and no lifted sum can wrap around that cyclic group.
The normalizing denominator is a positive standard constant multiple of
the count of such cubes. Thus the expansion gives a coordinate of (17)
with a positive standard average, up to a fixed positive normalizing factor,
on the convex polytope
\[
 \Omega=\{(b,b',h,n): b,b'\in I_2,\quad
       1\leq n+\omega\cdot h\leq N\text{ for all }\omega\in\{0,1\}^m\}.
                                                               \tag{19}
\]
In particular $\Omega\subseteq I_8^{s+1}$, and
$\left|\E_{x\in I_8^{s+1}}\one_\Omega(x)\sigma(x)\right|\gg1$.
All defining inequalities in (19) are affine linear; their number depends
only on $s$. The selection of biased pairs is used only for a lower bound
on the nonnegative norm average, not as an extra arbitrary cutoff in (19).

Published E.11 asks for dense subprogressions of a positive interval.
Partition $I_8^{s+1}$ into its finitely many sign orthants and coordinate-zero
hyperplanes. Since $\sigma$ is bounded, the latter have normalized mass
$O_s(1/N)$ and contribute $o(1)$. By the triangle inequality one open
orthant has positive standard weighted average. For its signs
$\varepsilon_i\in\{-1,1\}$, reflect $x_i=\varepsilon_i y_i$.
The reflected polytope is convex and the reflected nilcharacter
$\sigma_\varepsilon(y)=\sigma(\varepsilon y)$ has degree at most $s$.
Take every progression in E.11 to be the full positive interval $[8N]$.
The hypothesis is plain bias, hence bias of degree strictly below $s$.
E.11, with dimension $s+1$, degree $s\geq3$ and scale $8N$, gives
$\sigma_\varepsilon\in\Nil_{<s}(I_8^{s+1})$. Reflecting back gives
\[
 [\sigma]_{\Symb_s(I_8^{s+1})}=0.                        \tag{20}
\]
Only one orthant was used to obtain bias; E.11 itself supplies the conclusion
on the full symmetric box. No extrapolation from positive points is implicit.

Pull (18),(20) back along the fixed integer linear map
\[
 (h,n)\longmapsto(b,b',h_1,\ldots,h_m,n_{\rm aux})
                 =(h,n,n,\ldots,n,0).
\]
It maps $I_8^2$ into $I_8^{s+1}$. Lower-degree nilsequences pull back with
the same degree, so
\[
 s!\bigl(T(h;n^r)-T(n;h,n^{r-1})\bigr)=0
                    \quad\hbox{in }\Symb_s(I_8^2).       \tag{21}
\]
The sign in (18) can be removed because both sides equal zero. Apply E.13
on this symmetric box, with dimension two, degree $s\geq3$ and the fixed
positive integer $s!$. Its torsion-free conclusion gives
\[
 T(h;n^r)=T(n;h,n^{r-1})\quad\hbox{in }\Symb_s(I_8^2).  \tag{22}
\]
We restrict to $I_1\times J$ only after this use. E.13 is not being applied
directly to an asymmetric rectangle.

\section{Exact finite-component factorization}
Set $\Theta(n)=\widetilde\chi(n,\ldots,n)$, the same object as in LP05.
It is a unit degree-at-most-$s$ nilcharacter. Expand
$[\Theta(n+h)]-[\Theta(n)]$ by LP13. Terms with exactly one $h$ equal
\[
 T(h;n^r)+rT(n;h,n^{r-1})=sT(h;n^r)=[\chi(h,n)]          \tag{23}
\]
on $I_1\times J$, by (22),(3). Let $R(h,n)$ be the tensor product of
the actual pullbacks of $\widetilde\chi$ for all subsets of the $s$ slots
of cardinality at least two, putting $h$ in those slots and $n$ elsewhere.
There is no permutation of the distinguished first slot in this definition.
Every factor has multidegree at most $(i,s-i)$ for its subset size $i$:
coarsen the cube filtration by counting the coordinates assigned to $h$
and $n$, using all weights at least the index. The polynomial orbit,
rationality and top frequency are as in the LP05 block coarsening.
Their product is therefore supported on the union of these downsets,
contained in $\mathcal J$.

By the symbol equality just proved, the unit-vector expression
\[
 W(h,n)=\chi(h,n)\otimes\overline{\Theta(n+h)}\otimes\Theta(n)\otimes R(h,n)
                                                               \tag{24}
\]
is a degree-at-most-$(s-1)$ nilsequence on $I_1\times J$.
By the definition of the restricted nilsequence class, it has a global
representation whose restriction is (24); no global equality beyond this
rectangle is claimed. Such a degree filtration gives multidegree support
$i+j\leq s-1$ by assigning the group at $(i,j)$ its total-degree group.
Choose that representation and set $\Psi=\overline R\otimes W$.
It has support in (1) by the product filtration and is uniformly bounded.

Finally tensor (24) with
$\Theta(n+h)\otimes\overline{\Theta(n)}\otimes\overline R$.
Contract the matching pairs using their unit norms. Each contraction is
a fixed linear map between fixed finite tensor spaces. The result is exactly
$\chi$, so (2) holds component by component with fixed bounded coefficients.
This does not divide by any scalar coordinate or replace symbol equality
by an unproved scalar equality. It proves Theorem~\ref{thm:target}.

\section{Bounds, source differences and review boundary}
All finite coordinate selections, tensor dimensions, permutation counts
and parameter boxes depend only on fixed degree and geometric data.
Appendix A selection/uniformity supplies the finite family of lower tests
before the averaging manipulations. The three Cauchy--Schwarz steps use
common standard bounds, and each pigeonhole retains a standard positive
amount. Fixed parameter values may be unbounded limit integers; the project
shift theorem supplies uniform constructions for exactly that case.
All algebraic identities are finite and transfer coordinatewise, with the
all-limit vertical convention retained. E.11 and E.13 have the exact degree,
dimension, progression and symmetric-box domains stated above.

The proof retains, rather than erases, the published factorial history.
Theorem 7.4 on both pp.~1271 and 1315 prints the stricter hypothesis
``$<s-2$-correlates'', whereas the preceding Theorem 7.1 supplies
degree-at-most-$(s-2)$ correlation. The target here deliberately assumes
the latter, and Section 2 proves the required selection and shift step
under that weaker hypothesis. The printed Theorem 7.4 is not a premise.
In addition to the repeated shift index and alternating-derivative sign in
(13.14)--(13.17), the first two slots in printed Lemma 13.4 are interchanged
relative to (13.3),(13.10). This changes the sign in (13.13). Formula (18)
is the resulting coefficient for the actual unconjugated objects. The
published simplification in Lemma 13.3 also retains an $a$ where cancellation
requires $b$; its tail has degree at most $s-3$ in $n$. The scalar lower-factor
bookkeeping and the $I_6$ substitution above give their own exact formulas.
These source discrepancies are not counterexamples to the final zero relation.
Their required related uses and current representations must be checked
independently before a complete correction claim.

This packet proves a local composition relative to its stated project and
published premises. It does not reprove E.11/E.13's equidistribution inputs,
Proposition 5.6's published converse sources, Proposition 8.4, the frequency
analysis, or the final finite-family inverse-theorem assembly. Those are
direct dependencies or separate consumers, not silently discharged here.
\end{document}
